% !TEX root = ../main.tex

\chapter{Introduzione}

\section{Contesto generale}

Negli ultimi anni, la diffusione delle \textit{New Psychoactive Substances} (NPS) ha rappresentato una delle sfide più complesse per le forze dell’ordine e le agenzie doganali europee.
Le NPS sono molecole sintetiche progettate per imitare gli effetti di sostanze stupefacenti già note, ma con una struttura chimica modificata che le rende difficili da rilevare attraverso i metodi tradizionali. La loro rapida evoluzione, unita alla facilità di produzione in laboratori clandestini di piccole dimensioni, determina un contesto operativo caratterizzato da elevata incertezza e da un continuo aggiornamento dei repertori analitici.

Dal punto di vista chimico e operativo, una caratteristica distintiva delle NPS è la loro capacità di diffondersi nell’ambiente in forma di tracce volatili o aerosol durante le fasi di produzione, manipolazione e confezionamento.
Tali emissioni, pur avvenendo a concentrazioni estremamente basse, possono costituire un indicatore rilevante della presenza di attività illecite, soprattutto in prossimità di laboratori clandestini o siti di trasformazione temporanei.

A differenza delle sostanze tradizionali, la produzione di NPS è spesso associata a laboratori di piccole dimensioni, facilmente riconfigurabili e potenzialmente mobili, che operano in ambienti urbani o peri-urbani.
Questo rende inefficaci molte delle strategie di contrasto basate su controlli statici o su ispezioni mirate esclusivamente ai punti di distribuzione finale, spostando il focus delle operazioni investigative verso il rilevamento precoce delle emissioni ambientali.

La possibilità di intercettare segnali chimici dispersi nell’aria rappresenta, pertanto, un cambio di paradigma nel contrasto alla produzione di NPS, consentendo alle forze dell’ordine di individuare potenziali sorgenti prima che le sostanze entrino nella catena di distribuzione.

Queste caratteristiche rendono evidente come il problema del rilevamento delle NPS non possa essere affrontato esclusivamente attraverso analisi di laboratorio o controlli a posteriori, ma richieda strumenti capaci di operare direttamente sul campo, in tempo reale e in condizioni ambientali non controllate.

In scenari reali, gli operatori devono spesso intervenire in ambienti dinamici e poco strutturati, come aree urbane densamente popolate o zone industriali dismesse, dove la presenza di sostanze volatili o residui chimici può costituire un indicatore cruciale di attività illecite.
In questo contesto, la possibilità di impiegare sistemi mobili in grado di rilevare, analizzare e interpretare in tempo reale la presenza di composti chimici rappresenta un obiettivo di grande rilevanza sia scientifica sia operativa.

A livello europeo, il fenomeno delle sostanze sintetiche rappresenta una minaccia crescente sia per la sicurezza pubblica sia per la salute dei cittadini. Negli ultimi anni, organizzazioni criminali distribuite in diversi Paesi membri hanno adottato modelli produttivi estremamente flessibili, fondati su piccoli laboratori mobili, precursori facilmente reperibili e tecniche chimiche adattive. Questa dinamica rende la produzione di NPS difficilmente individuabile e consente una rapida introduzione di nuove varianti, spesso non ancora incluse nei repertori analitici ufficiali delle agenzie governative.

Tale scenario è ulteriormente aggravato dal fatto che molte attività di produzione e distribuzione avvengono in aree urbane complesse, caratterizzate da elevata densità abitativa e da infrastrutture fisiche che ostacolano sia le operazioni di rilevamento sia la modellazione della dispersione di sostanze volatili. Le forze dell’ordine europee hanno quindi espresso la necessità di strumenti capaci di supportare il processo decisionale sul campo, fornendo analisi affidabili, contestualizzate e rapidamente verificabili.

In risposta a queste esigenze, il progetto \textsc{PENTION} si propone di sviluppare una piattaforma digitale e fisica orientata al monitoraggio, alla modellazione e al rilevamento precoce di attività legate alla sintesi di droghe sintetiche. L’iniziativa mira a combinare modelli predittivi, tecniche di analisi automatizzata, simulazioni fisiche e strumenti di supporto investigativo, contribuendo a colmare il gap tecnologico attualmente presente nelle operazioni delle LEA europee.

Tuttavia, l’identificazione automatica delle NPS presenta diversi ostacoli: gli spettri di massa ottenuti sul campo possono essere rumorosi, parziali o deteriorati; le condizioni ambientali influenzano fortemente la dispersione chimica; e molti strumenti attuali non sono pensati per fornire risultati \textit{tracciabili} o utilizzabili in un contesto investigativo.
È dunque necessario progettare sistemi che integrino modelli fisici, tecniche di machine learning, simulazioni ambientali e meccanismi di audit forense, in modo da garantire robustezza, spiegabilità e affidabilità del processo decisionale.

Uno degli aspetti critici evidenziati dagli operatori del settore riguarda la mancanza di sistemi integrati in grado di combinare misurazioni ambientali, modelli fisici e algoritmi di intelligenza artificiale all'interno di un'unica piattaforma operativa. Le soluzioni attualmente disponibili sono spesso frammentate, non interoperabili e prive di un meccanismo strutturato di tracciabilità che consenta la validazione forense dei risultati. Inoltre, la maggior parte degli strumenti commerciali non è progettata per un utilizzo mobile né per l’analisi contestuale in aree urbane reali, limitando notevolmente la capacità delle agenzie di condurre operazioni dinamiche e tempestive.

La necessità di superare tali limitazioni richiede la progettazione di un sistema capace di integrare diverse fonti informative, aggiornarsi dinamicamente durante la perlustrazione del territorio e produrre evidenze digitali verificabili. È proprio su questa esigenza che si fonda la visione del progetto \textsc{PENTION}, orientato alla creazione di una piattaforma cyber-fisica avanzata per il rilevamento e l’interpretazione di fenomeni chimici riconducibili alla produzione di sostanze sintetiche.

\section{PENTION-S e PENTION-M: due sistemi complementari}

Il progetto \textsc{PENTION} prevede lo sviluppo coordinato di due piattaforme autonome e complementari, ciascuna progettata per rispondere a esigenze operative specifiche delle forze dell’ordine europee.
\textsc{PENTION-S} è orientato al monitoraggio statico e continuo attraverso l’installazione di sensori fissi in infrastrutture critiche quali aeroporti, centri logistici, uffici postali e aree ad accesso controllato. In tali contesti, la possibilità di acquisire misurazioni da sorgenti posizionate stabilmente consente analisi regolari, accumulo di evidenze nel tempo e rilevamento precoce di anomalie chimiche in ambienti generalmente strutturati e prevedibili.

\textsc{PENTION-M}, invece, nasce per affrontare scenari radicalmente diversi, caratterizzati da mobilità, variabilità ambientale e necessità di intervento rapido.
Le attività di perlustrazione urbana, i controlli in prossimità di porti e confini, le investigazioni in aree industriali o rurali e, più in generale, tutte le situazioni in cui gli operatori devono spostarsi attivamente sul territorio richiedono una piattaforma capace di raccogliere, elaborare e interpretare dati in tempo reale durante il movimento. In questi scenari, la mobilità non è un'opzione accessoria, ma un requisito essenziale.

Pur condividendo alcune finalità generali del progetto \textsc{PENTION} — in particolare il rilevamento precoce di sostanze chimiche e il supporto data-driven alle decisioni operative — i due sistemi non vanno intesi come versioni successive dello stesso artefatto software. Essi rappresentano due risposte progettuali distinte a problemi differenti: stabilità e continuità nel caso di \textsc{PENTION-S}, flessibilità e reattività nel caso di \textsc{PENTION-M}.

Nel contesto di questa tesi, l’analisi iniziale di \textsc{PENTION-S} ha svolto un ruolo conoscitivo importante per comprendere le logiche di modellazione, le assunzioni fisiche e le scelte architetturali alla base del sistema statico. Tuttavia, \textsc{PENTION-M} è stato progettato come un artefatto indipendente, dotato di un’architettura dedicata, di modelli specifici per la mobilità e di una pipeline operativa completamente ripensata per l’elaborazione di dati raccolti durante il movimento del veicolo.

\section{Obiettivi della tesi}

L'obiettivo principale di questa tesi è la progettazione e la realizzazione di un sistema in grado di simulare, in maniera realistica e operativamente coerente, il comportamento di un laboratorio mobile dedicato al rilevamento e all’analisi di sostanze psicoattive emergenti.
La piattaforma risultante, denominata \textsc{PENTION-M}, è stata progettata come sistema autonomo dedicato al rilevamento mobile, distinto ma complementare rispetto a \textsc{PENTION-S}, con cui condivide solo alcuni principi concettuali individuati nelle fasi preliminari dello studio.

A partire da questo obiettivo generale, il lavoro può essere articolato in una serie di obiettivi specifici:

\begin{itemize}
	\item \textbf{Progettare un modello di simulazione fisica realistico} della dispersione di sostanze chimiche in ambiente urbano, capace di adattarsi a condizioni meteo variabili e alla presenza di edifici, integrando una rappresentazione dinamica del movimento del veicolo.
	\item \textbf{Integrare modelli \textit{Physics-Informed Machine Learning} (PIML)} per la correzione e l’interpretazione delle mappe di dispersione, garantendo coerenza con i vincoli fisici del dominio operativo.
	\item \textbf{Sviluppare un modello di localizzazione della sorgente} basato su informazioni fisiche derivate, in grado di stimare la posizione del rilascio con stabilità e robustezza.
	\item \textbf{Costruire un classificatore di sostanze NPS affidabile}, facendo uso di un dataset ampio e derivato da molteplici fonti, e dotato di un meccanismo di calibrazione probabilistica adattiva per la riduzione dell’overconfidence.
	\item \textbf{Progettare una pipeline MLOps completa}, comprendente ingestion, validazione, calcolo delle feature fisiche, inferenza dei modelli, monitoring dei parametri operativi (latenza, drift, versione del modello, confidenza) e strategie per il retraining.
	\item \textbf{Realizzare un sistema di logging forense e chain of custody}, tramite la generazione di un bundle di evidenze digitali contenente dati, modelli, parametri e firme crittografiche, con l'obiettivo di garantire tracciabilità, integrità e verificabilità dell'intero processo.
	\item \textbf{Sviluppare un'interfaccia interattiva in tempo reale} che consenta di orchestrare la simulazione, visualizzare i risultati e riprodurre fedelmente il flusso operativo di un laboratorio mobile.
\end{itemize}

Questi obiettivi si collocano in continuità con le finalità più ampie del progetto \textsc{PENTION}, che include tra le sue linee strategiche la riduzione della produzione di droghe sintetiche in Europa, il potenziamento delle capacità di identificazione precoce di precursori e sottoprodotti chimici, e il supporto alle agenzie di sicurezza attraverso strumenti affidabili, spiegabili e legalmente ammissibili. Il sistema proposto in questa tesi contribuisce a tali finalità fornendo un prototipo operativo che integra in modo coerente simulazione fisica, analisi predittiva e generazione di evidenze digitali verificabili.

Nel loro insieme, questi obiettivi concorrono alla costruzione di un prototipo avanzato che unisce aspetti di simulazione fisica, apprendimento automatico, ingegneria del software, modellazione operativa e requisiti forensi.
La tesi si colloca dunque all'intersezione tra sistemi cyber-fisici, intelligenza artificiale applicata e tecnologie di supporto alle attività investigative, con l'intento di offrire un contributo concreto allo sviluppo di piattaforme intelligenti per la sicurezza pubblica.

\section{Metodologia adottata}

La natura del problema affrontato in questa tesi, che coinvolge aspetti di simulazione fisica, modellazione predittiva, progettazione software e requisiti operativi provenienti da contesti investigativi reali, richiede un approccio metodologico capace di integrare attività di analisi, progettazione, implementazione e valutazione in un ciclo iterativo e guidato dagli stakeholder.
Per questo motivo il lavoro si colloca all'interno del paradigma della \textit{Design Science Research} (DSR), una metodologia ampiamente adottata nella ricerca ingegneristica e informatica per la costruzione di artefatti innovativi orientati alla soluzione di problemi concreti.

La DSR pone al centro del processo la progettazione di un artefatto -- nel caso di questa tesi, il sistema \textsc{PENTION-M} -- inteso come combinazione strutturata di modelli fisici, algoritmi di machine learning, pipeline operative e componenti software.
L'efficacia dell’artefatto non deriva unicamente dalla sua implementazione tecnica, ma dalla sua capacità di rispondere ai requisiti emersi da un contesto reale, e dalla possibilità di valutarlo attraverso iterazioni successive che riflettono feedback, vincoli e obiettivi operativi.

In questo lavoro, la metodologia DSR si articola in tre dimensioni principali:

\begin{itemize}
	\item \textbf{Relevance}, ovvero la rilevanza del problema: il rilevamento di sostanze NPS richiede strumenti robusti, tracciabili e operativamente affidabili, come confermato dalle attività di Requirements Engineering condotte con le forze dell’ordine olandesi e con i partner del progetto;
	\item \textbf{Rigor}, ovvero il rigore scientifico: il sistema è fondato su basi teoriche consolidate, quali modelli di dispersione gaussiana, metodologie PIML, tecniche di calibrazione probabilistica e principi di MLOps per l’analisi forense e la gestione del ciclo di vita dei modelli;
	\item \textbf{Design \& Evaluation}, ossia progettazione e valutazione iterativa: l’artefatto è stato sviluppato attraverso versioni incrementali, ciascuna valutata tramite feedback degli stakeholder, analisi dei risultati e verifica della coerenza fisica e funzionale.
\end{itemize}

Un ulteriore elemento distintivo del processo adottato riguarda la forte interdisciplinarità del progetto, che combina competenze di chimica analitica, modellazione atmosferica, data science, ingegneria del software e requisiti provenienti direttamente dal dominio investigativo. Tale integrazione riflette pienamente la visione espressa nella proposta progettuale, secondo cui l’efficacia di strumenti come \textsc{PENTION-M} dipende dalla capacità di tradurre esigenze operative complesse in artefatti ingegneristici robusti e scientificamente fondati.

L’adozione del paradigma DSR consente di strutturare la tesi non solo come un lavoro di implementazione tecnica, ma come un processo di ricerca che parte dall’identificazione del problema, passa attraverso la definizione degli obiettivi e la costruzione di una soluzione, e si conclude con una valutazione integrata nel ciclo stesso di sviluppo.
Nel Capitolo~\ref{chap:dsr} verrà presentata in dettaglio l’applicazione concreta del Design Science Research al progetto \textsc{PENTION-M}, includendo il ciclo iterativo, il ruolo degli stakeholder e l’impatto dei requisiti sul design finale del sistema.

\section{Struttura della tesi}

La presente tesi è articolata in sette capitoli, ciascuno dei quali affronta un aspetto specifico del percorso di progettazione, sviluppo e valutazione del sistema \textsc{PENTION-M}.
Di seguito si fornisce una panoramica della struttura complessiva dell’elaborato.

\begin{itemize}
	\item \textbf{Capitolo 1 -- Introduzione}: introduce il contesto operativo e scientifico del progetto, chiarisce il ruolo complementare dei sistemi \textsc{PENTION-S} e \textsc{PENTION-M}, definisce gli obiettivi della tesi e presenta la metodologia adottata.
	\item \textbf{Capitolo 2 -- Background, Stato dell’Arte e Fondamenti}: offre una rassegna dei concetti teorici e tecnologici necessari per comprendere il sistema sviluppato, includendo una panoramica sulle NPS, sugli spettri EI, sulle principali banche dati utilizzate, sui modelli di dispersione, sulle metodologie PIML, sulle pratiche MLOps e sull’architettura originale di \textsc{PENTION-S}.
	\item \textbf{Capitolo 3 -- Metodologia Design Science Research}: descrive l’approccio metodologico adottato, illustra il processo di \textit{Requirements Engineering}, presenta il Data Flow Model delle iterazioni progettuali e documenta il ruolo degli stakeholder nel ciclo di sviluppo.
	\item \textbf{Capitolo 4 -- Architettura e Realizzazione di \textsc{PENTION-M}}: costituisce il nucleo tecnico della tesi, mostrando in dettaglio i moduli software, i modelli fisici, i modelli PIML, la pipeline MLOps, il classificatore di sostanze, il sistema di logging forense e l’interfaccia operativa.
	\item \textbf{Capitolo 5 -- Validazione e Discussione dei Risultati}: presenta la valutazione dell’artefatto secondo il paradigma DSR, analizza il contributo degli stakeholder, valuta prestazioni e robustezza dei modelli, e include un \textit{Experience Report} con le principali lezioni apprese.
	\item \textbf{Capitolo 6 -- Limitazioni e Sviluppi Futuri}: discute le attuali limitazioni tecniche e operative del sistema e suggerisce possibili direzioni di miglioramento, tra cui la conduzione di un field study e l’integrazione con hardware reale.
	\item \textbf{Capitolo 7 -- Conclusioni}: sintetizza i contributi principali del lavoro, evidenzia l’impatto scientifico e operativo del sistema proposto e riflette sul potenziale di \textsc{PENTION-M} come strumento innovativo per il supporto investigativo.
\end{itemize}

Questa struttura riflette la progressione logica del lavoro svolto, permettendo al lettore di seguire in modo chiaro il percorso che va dall’analisi del problema alla progettazione dell’artefatto, fino alla sua valutazione e contestualizzazione nei possibili scenari futuri.